\paragraph{Computational consequences.}
Let $G$ be the number of active destination groups, $N$ the number of graph
nodes, and $A$ the number of links. Routing preparation requires one backward
pass and probability construction per group, while each subsequent cached
evaluation requires only the forward loading passes. The principal incremental
cache arrays scale as $O(GA)$ for probabilities and masks. Compact OD grouping
can reduce both $G$ and this storage when frozen-zero demand removes complete
destination groups.

Table~\ref{tab:fixed-routing-benchmarks} reports illustrative CPU measurements
from the package examples. Timings are warm medians using 32-bit arithmetic and
are machine-dependent; numerical equivalence is tested separately and these
values are not used as test thresholds.

\begin{table}[htbp]
\centering
\begin{tabular}{lrrrrr}
\hline
Example and layout & $G$ & Cache (MiB) & Forward & Value--gradient & Speedup \\
\hline
Simple example, compact & 7 & 0.20 & $9.44\!\to\!2.75$ ms & $12.78\!\to\!6.21$ ms & $2.06\times$ \\
Geneva, compact & 18 & 3.43 & $0.449\!\to\!0.109$ s & $0.537\!\to\!0.180$ s & $2.98\times$ \\
Geneva, full & 62 & 11.83 & $1.552\!\to\!0.383$ s & $1.788\!\to\!0.591$ s & $3.03\times$ \\
\hline
\end{tabular}
\caption{Dynamic-to-cached fixed-routing measurements. The final column reports
the likelihood value-and-gradient speedup; forward speedups range from
$3.43\times$ to $4.12\times$.}
\label{tab:fixed-routing-benchmarks}
\end{table}

